\documentclass[11pt,a4paper]{article}
\usepackage[T1]{fontenc}
\usepackage{lmodern,amsmath,amssymb,amsthm,mathtools,booktabs,microtype}
\usepackage[margin=25mm]{geometry}
\usepackage[colorlinks=true,linkcolor=blue,urlcolor=blue,citecolor=blue]{hyperref}
\newtheorem{theorem}{Theorem}[section]
\newtheorem{lemma}[theorem]{Lemma}
\newtheorem{proposition}[theorem]{Proposition}
\newtheorem{corollary}[theorem]{Corollary}
\newcommand{\Q}{\mathbb Q}\newcommand{\Z}{\mathbb Z}
\DeclareMathOperator{\Disc}{Disc}\DeclareMathOperator{\supp}{supp}
\setlength{\emergencystretch}{2em}
\hypersetup{pdftitle={Separating Quadratic Roots in Finite Algebras and Smaller Degrees for Diophantine Presentations}}
\title{Separating Quadratic Roots in Finite Algebras\\and Smaller Degrees for Diophantine Presentations}
\author{}\date{13 November 2026}
\begin{document}\maketitle
\begin{abstract}
We give an explicit polynomial encoding a rational point of a finite algebra
subject to a nonzero condition. A triangular linear combination separates
quadratic roots without a lower bound on their nonzero radicands. Separate
bounds for the base roots and the radicands keep the polynomial weights
small. Characteristic polynomials are truncated to the dimension of the
relevant nonzero part. The construction includes vanishing radicands,
nonreal roots, repeated characteristic values, and nonreduced excluded
components. Applied to the existing arithmetic sources, it gives degree
bounds below $2^{78}$ in 25 variables, below $2^{80}$ in the fully specified
27-variable construction, and below $2^{38}$ in both conditional
13-variable constructions. The fully specified coefficient-height bound
improves to $2^{2^{92}}$. The variable counts, arithmetic hypotheses and
sufficient coefficient parameters are unchanged.
\end{abstract}

\section{Applications and their precise scope}
Every variable count includes the distinguished input. The degree is total
degree, and coefficient height is the largest absolute integer coefficient.
The following are bounds for new polynomials with exactly the same rational
input projections as the source constructions in the repository.
\begin{center}\small
\begin{tabular}{lrrr}
\toprule
Source & Variables & Degree at most & Support less than\\\midrule
Inherited arithmetic premises &25&185517288563368898320200&$2^{1815}$\\
Fully specified, inherited premises &27&679086760181946108087584&$2^{2009}$\\
Weak $abc$ at $121/100$ &13&207969001841&$2^{437}$\\
Weak $abc$ at $16/15$ &13&199368913265&$2^{436}$\\\bottomrule
\end{tabular}\end{center}
The last two rows also retain the elliptic group and height, pole-realization,
and prime-pattern inputs of their sources. The first two retain their
additional modular and contact-height inputs. This paper changes none of
those arithmetic statements. It proves a new algebraic conversion of the
same sources. It is not an independent verification of their arithmetic
foundations.

For 25 variables the sufficient odd coefficient parameter remains ineffective.
For 27 variables the existing certified value $L=20483$ remains sufficient
under the same premises. In both 13-variable families the sufficient odd
shift $B$ remains ineffective; the numerical $B=3$ examples are not
certified defining members. Multiplication by the input includes zero in
the two 13-variable constructions, exactly as in their sources.

The source polynomials have not been expanded. Thus the displayed degrees
and support sizes are upper bounds, not exact degrees or counts. The
coefficient specification in Section~\ref{sec:formula} consists entirely of
finite integer-polynomial operations.

\section{Finite algebras and the nonzero condition}
Let $R=\Z[x_1,\ldots,x_a]$. Let
\[
 B=R[z_1,\ldots,z_p]/(f_1(z_1),\ldots,f_p(z_p)),\qquad
 A=B[q_1,\ldots,q_s]/(q_j^2-c_j)_{j=1}^s,
\]
where $p,s\ge1$, the $f_i\in R[Z]$ are monic of degrees $d_i\ge2$,
and $c_j\in B$. Suppose each $f_i$ is separable at every rational
specialization of the base coordinates. Put
\[
 b=\prod_i d_i,\qquad n=2^s b.
\]
The usual monomials of exponent below the defining degrees are bases of
$B$ and $A$. Write $M_v^B$ or $M_v^A$ for multiplication by $v$ in these
bases. Fix $g\in B$ and set
\[
 G=gq_1\cdots q_s,\qquad T=g^2c_1\cdots c_s\in B.
\]
At a rational base point our target condition is
\begin{equation}\label{eq:target}
 \text{there exists a $\Q$-algebra map $A\longrightarrow\Q$ with $G\ne0$.}
\end{equation}
There are no additional residual equations. The existing source certificates
verify this form of the nonzero condition and the independence of each
$c_j$ from the $q$-generators.

At a rational base point $B$ is an etale algebra. On its factors where
$T\ne0$, all the $c_j$ and $g$ are nonzero, and the corresponding part of
$A$ is also etale. The other factors of $A$ may have nilpotents. They will
not be discarded by an unjustified division.

\section{A triangular combination of quadratic roots}
\begin{lemma}\label{lem:triangular}
Let $S\ge1$ be real and let $c_1,\ldots,c_s\in\mathbb C^\times$ satisfy
$|c_j|\le S$. For each choice $q_j^2=c_j$, put
\[
 \phi(q_1,\ldots,q_s)
 =\frac1{2S}\sum_{j=1}^s(4S)^{1-j}
       \left(\prod_{k<j}c_k\right)q_j.
\]
Its values on the $2^s$ choices of roots are distinct and have absolute
value at most $2/3$. The same absolute-value bound holds when some
$c_j$ vanish, with the corresponding root equal to zero.
\end{lemma}
\begin{proof}
Every $|q_j|\le\sqrt S$. The sum of the absolute values is at most
\[
 \frac{\sqrt S}{2S}\sum_{j=1}^s4^{1-j}
 <\frac{2}{3\sqrt S}\le\frac23.
\]
Consider two distinct choices, and let $j$ be their first differing index.
The common earlier factor is nonzero. After dividing the difference by
that factor, its leading term is $q_j-\widetilde q_j=2q_j$. For the tail,
use $|c_j|=|q_j|^2$ and $|q_k|\le\sqrt S$. The ratio of its absolute
value to $2|q_j|$ is bounded above by
\[
 \frac{|q_j|}{\sqrt S}\sum_{\ell\ge1}4^{-\ell}\le\frac13.
\]
The leading term cannot be cancelled. This proof uses no lower bound for
$|c_j|$ and applies equally to nonreal roots.
\end{proof}

Write $f_i(Z)=Z^{d_i}+\sum_{k<d_i}f_{ik}Z^k$ and define scalar polynomials
\begin{equation}\label{eq:bounds}
 Z_0=2+\sum_{i,k<d_i}f_{ik}^{\,2},\qquad
 S=b\left(1+\sum_{j=1}^s\sum_{u,v=1}^b(M_{c_j}^B)_{uv}^{\,2}\right).
\end{equation}
At every rational base point $Z_0\ge2$ bounds the absolute values of the
point roots. Indeed Cauchy's monic root bound is $1+\max_k|f_{ik}|$,
and $1+t\le2+t^2$ for $t\ge0$. Also $S\ge b\ge1$ bounds the absolute
values of all character values of $c_j$: the spectral radius is at most
$b$ times the largest absolute entry, and $|t|\le1+t^2$.

Set
\begin{gather*}
 \delta_i=\Disc(f_i),\qquad
 \kappa=\max_i\frac{d_i(d_i-1)-2}{2},\qquad
 D=\prod_i\delta_i^2,\\
 K=(2Z_0)^\kappa\prod_i(1+\delta_i^2),\qquad C=D+4Z_0K.
\end{gather*}
At a rational base point $D,K>0$. Put $\epsilon=D/K$ and $L=C/D$.
Then $0<\epsilon\le1$ and $L=1+4Z_0/\epsilon$.

\begin{lemma}\label{lem:pointsep}
Any two distinct roots of any $f_i$ differ in absolute value by at least
$\epsilon$.
\end{lemma}
\begin{proof}
If $\Delta$ is the chosen root difference, the discriminant product gives
\[
 |\delta_i|\le |\Delta|^2(2Z_0)^{d_i(d_i-1)-2}.
\]
For $t>0$, $t^2/(1+t^2)\le\sqrt t$. Therefore
\[
 \epsilon\le \frac{\delta_i^2}{1+\delta_i^2}(2Z_0)^{-\kappa}
 \le\frac{\sqrt{|\delta_i|}}
 {(2Z_0)^{(d_i(d_i-1)-2)/2}}\le|\Delta|.
\]
The exponent is the maximum needed for a single polynomial, not the sum
of all discriminant exponents.
\end{proof}

Define in $A$
\begin{gather*}
 J=4S,\qquad Q=2S J^{s-1},\qquad
 \Xi=\sum_{j=1}^sJ^{s-j}\left(\prod_{k<j}c_k\right)q_j,\\
 \Omega=Q\left(C^{p+1}+\sum_{i=1}^pD^iC^{p+1-i}z_i\right)+D^{p+1}\Xi.
\end{gather*}
These expressions are polynomial and involve no division.

\begin{proposition}\label{prop:sep}
At every rational base point, $\Omega$ has distinct nonzero values on all
geometric characters of $A$ with $G\ne0$.
\end{proposition}
\begin{proof}
Divide its values by the nonzero rational number $QD^{p+1}$. The result is
\[
 \alpha=L^{p+1}+\sum_{i=1}^pL^{p+1-i}z_i+\phi.
\]
For two characters first differing at $z_i$, put $k=p+1-i\ge1$.
The leading difference has magnitude at least $\epsilon L^k$.
The remaining difference is strictly less than
\[
 2Z_0\sum_{\ell=0}^{k-1}L^\ell
 <\frac{2Z_0}{L-1}L^k=\frac\epsilon2L^k.
\]
Here the difference of the two $\phi$ values is at most $4/3\le2Z_0$.
If all point generators agree, the radicands agree and are nonzero;
Lemma~\ref{lem:triangular} separates the remaining root choices.
Finally
\[
 |\alpha-L^{p+1}|\le Z_0\sum_{\ell=0}^pL^\ell
 <\frac{Z_0}{L-1}L^{p+1}\le\frac14L^{p+1}.
\]
Thus no value is zero. Notice that this last nonvanishing estimate also
holds on excluded characters, although their separation is unnecessary.
\end{proof}

\section{A single polynomial and all degree-drop cases}\label{sec:formula}
Define scalar coefficients by
\[
 \det_B(I-vM_T)=\sum_{j=0}^b a_jv^j,\qquad a_0=1.
\]
For each $1\le r\le b$, set
\[
 h_r(X)=\sum_{j=0}^r a_jX^{r-j},\qquad E_r=a_r-h_r(T),\qquad
 \omega_r=E_r\Omega,\qquad d_r=2^s r.
\]
Let $H_r(U)$ be the truncation through degree $d_r$ of
$\det_A(I-UM_{\omega_r})$. Equivalently, it is the sum of its principal-minor
terms with sizes from zero through $d_r$. Finally set
\begin{equation}\label{eq:polynomial}
 P=H_b\prod_{r=1}^{b-1}\left(H_r^2+\sum_{j=r+1}^b a_j^2\right).
\end{equation}
This is an explicit integer polynomial in the $a$ retained base coordinates
and just one further rational coordinate $U$.

\begin{theorem}\label{thm:projection}
For every rational base point, condition~\eqref{eq:target} holds if and
only if $P(U)=0$ for some $U\in\Q$.
\end{theorem}
\begin{proof}
Let $r_0$ be the number of geometric characters of $B$ on which $T\ne0$,
counted individually, not by distinct values of $T$. Because $B$ is etale,
\[
 \det_B(I-vM_T)=\prod_\chi(1-vT(\chi))
\]
has degree exactly $r_0$. If $r=r_0>0$, then $a_r\ne0$ and
$h_r(X)$ is the product of $X-T(\chi)$ over the $r$ nonzero character
values, with repetitions retained. Hence $E_r=a_r$ on this part and
$E_r=0$ on $T=0$.

If $r>r_0$, then $a_r=0$ and
$h_r(X)=X^{r-r_0}h_{r_0}(X)$, with $h_0=1$ when $r_0=0$.
This vanishes at every value of $T$. Thus $E_r=0$ in $B$, and also in $A$;
$H_r=1$. This explains why no additional factor $a_r^2$ is required in
$\omega_r$. If $r<r_0$, a nonzero tail coefficient makes
$H_r^2+\sum_{j>r}a_j^2$ strictly positive at every rational $U$.
Consequently any rational zero of~\eqref{eq:polynomial} belongs to the
unique correct case $r=r_0>0$.

In that case $E_r$ kills the entire excluded part, including any nilpotents.
The remaining part is etale of dimension $d_r$. On it $\omega_r=a_r\Omega$
has distinct nonzero character values by Proposition~\ref{prop:sep}.
Thus the full determinant has degree $d_r$, so the truncation is exact.
Its rational zero gives a nonzero rational character value of $\omega_r$
of multiplicity one. The corresponding one-dimensional rational kernel
of $I-UM_{\omega_r}$ is preserved by every native multiplication matrix.
Their actions on this line are rational scalars satisfying the algebra
relations. They define a rational character on which $G\ne0$.

Conversely, a rational character with $G\ne0$ has a nonzero rational
value of $\omega_{r_0}$. Its reciprocal is a root of $H_{r_0}$, and all
higher $a_j$ vanish. This gives a zero of $P$. If $r_0=0$, every $H_r=1$
and $P$ has no zero. All assertions concern the complete rational zero
set, not a generic open subset.
\end{proof}

\section{Degree and support bounds}
Give the base coordinates weight one, the $z_i$ weights $w_i$, and the
$q_j$ weights $v_j$, so that monic reduction does not increase weight.
Suppose $\operatorname{wt}(G)\le w_G$. Put
\[
 t=2w_G,\quad u=\sum_i d_i(d_i-1)w_i,\quad
 b_w=\sum_i(d_i-1)w_i.
\]
The scalar bounds in~\eqref{eq:bounds} have degree at most
\[
 z_d=2\max_i d_iw_i,\qquad s_d=2\max_j(2v_j+b_w).
\]
Indeed an entry of multiplication by $c_j$ has degree at most
$2v_j+$ column weight $-$ row weight. Define
\[
 d_d=2u,\qquad c_d=(\kappa+1)z_d+2u,
\]
and
\begin{align*}
 A_0&=ss_d+\max\left\{(p+1)c_d,
        \max_{1\le i\le p}(id_d+(p+1-i)c_d+w_i)\right\},\\
 A_1&=(p+1)d_d+\max_{1\le j\le s}
        \left((s-j)s_d+2\sum_{k<j}v_k+v_j\right),\\
 A_*&=\max(A_0,A_1).
\end{align*}
These give $\operatorname{wt}(\Omega)\le A_*$.

\begin{theorem}\label{thm:degree}
Put $h_r=2^sr(rt+A_*+1)$. Then
\begin{align}
 \deg P&\le h_b+2\sum_{r=1}^{b-1}\max(h_r,bt)=:\mathcal D,\label{eq:D}\\
 \deg_U P&\le2^sb^2=:\mathcal E.\label{eq:E}
\end{align}
If $V=a+1$ is the total number of coordinates and $\mathcal E\le\mathcal D$,
then
\[
 \#\supp(P)\le
 \binom{\mathcal D+V}{V}-\binom{\mathcal D-\mathcal E+V-1}{V}.
\]
\end{theorem}
\begin{proof}
Principal minors cancel row and column basis weights. Thus
$\deg a_j\le jt$. The constant term cancels in $E_r=a_r-h_r(T)$,
leaving $r$ terms of weight at most $rt$. Hence $\omega_r$ has weight
at most $rt+A_*$. A principal minor of size $j$, including the factor
$U^j$, has total degree at most $j(rt+A_*+1)$. Only $j\le2^sr$
occur in $H_r$. Each tail coefficient has degree at most $bt$.
This proves~\eqref{eq:D}. The bound~\eqref{eq:E} follows from
$2^s(b+2\sum_{r<b}r)=2^sb^2$. Counting monomials by the exponent of
$U$ gives the displayed binomial difference.
\end{proof}
For independent numerical checking, the sum of the $h_r$ with the indicated
multiplicities is
\[
 2^s\left[t\frac{2b^3+b}{3}+(A_*+1)b^2\right].
\]
The correction for the initial interval $1\le r\le r_*$ where $h_r<bt$
is $2\sum_{r=1}^{r_*}(bt-h_r)$. The direct summation and this closed
formula are checked separately in the accompanying programs.

\section{Coefficient bounds}
We supply the bounds used for the fully specified 27-variable polynomial;
the same formulas also give correctly labelled bounds for the ineffective
families' numerical fixtures. Write $\ell(m)=\lceil\log_2m\rceil$ for
positive integers $m$. Use the weighted column norm over the coefficient
$\ell^1$ norm. Suppose
\[
 \|M_{z_i}\|\le2^{\gamma_i},\quad
 \|M_{q_j}\|\le2^{\nu_j},\quad
 \|M_G\|\le2^\eta,
\]
and $\log_2\|\delta_i\|_1\le\lambda_i$.
Put $b_\gamma=\sum_i(d_i-1)\gamma_i$ and $v=\sum_i\lambda_i$.
Valid scalar logarithmic norm bounds are
\begin{align*}
 z_l&=\ell(2+\sum_i d_i)+2\max_i d_i\gamma_i,\\
 s_l&=\ell(b)+\ell(1+sb^2)+2\max_j(2\nu_j+b_\gamma),\\
 d_l&=2v,\qquad
 c_l=\kappa(1+z_l)+p+d_l+z_l+3,\\
 q_l&=ss_l+2s-1,\\
 \xi_l&=\ell(s)+\max_j\left((s-j)(s_l+2)
                         +2\sum_{k<j}\nu_k+\nu_j\right).
\end{align*}
Here the restriction of multiplication by $q_j^2$ to $B$ bounds
$M_{c_j}^B$ in the inherited basis scaling. Hence these estimates do not
assume that the unweighted basis has operator norm one.

A bound for $\log_2\|M_\Omega\|$ is
\begin{align*}
 a_l=1+\max\bigg\{&q_l+\ell(p+1)+
       \max\big((p+1)c_l,\max_i(id_l+(p+1-i)c_l+\gamma_i)\big),\\
       &(p+1)d_l+\xi_l\bigg\}.
\end{align*}
Set $\mu=2\eta+\ell(b)$ and
\[
 k_r=2^sr\big(r\mu+\ell(r)+a_l+\ell(n)+1\big).
\]
Then
\begin{equation}\label{eq:norm}
 \log_2\|P\|_1\le k_b+
 \sum_{r=1}^{b-1}\left(2\max(k_r,b\mu)+\ell(b+1)\right).
\end{equation}
To verify this, principal-minor bounds give
$\|a_j\|_1\le\binom bj2^{2j\eta}\le2^{j\mu}$.
The $r$ terms of $E_r$ give norm at most $2^{r\mu+\ell(r)}$.
For the truncated determinant, each principal minor has norm bounded by
the corresponding column-norm product. Bounding $\binom nj\le n^j$
and $d_r+1\le2^{d_r}$ yields $k_r$. Squaring, adding at most $b+1$
terms, and multiplying the factors proves~\eqref{eq:norm}.

For the certified parameter $L=20483$ in the 27-variable source this gives
\[
 \log_2\|P\|_1\le4582336810897811998941387501<2^{92},
 \qquad H(P)<2^{2^{92}}.
\]
For the three numerical fixtures the respective logarithmic bounds are
\[
 \begin{array}{c|r}
 25\text{ variables},\ L=3&1154691810099707471428725\\
 13\text{ variables},\ B=3,\ \kappa=121/100&509178610649\\
 13\text{ variables},\ B=3,\ \kappa=16/15&496284465113.
 \end{array}
\]
None of these three fixture parameters is asserted to define the integers.

\section{Application checks and reproducibility}
The exact source versions are pinned in
\begin{center}\small
\nolinkurl{computation/few_variables/triangular_elimination_v24/inputs.json}.
\end{center}
The point counts are nine, nine, three and three; their scalar monic
polynomials all have degree four, weight six, and logarithmic operator
norm bound fourteen. The pole counts are two, three, two and two.
The source certificates include $G=g\prod q_j$, no additional residual
equations, and complete finite coefficient circuits.

For clarity, the source point polynomials have the form
$Y^4-4a'Y^2+8(b')^2$, obtained from $h^4-D_0=a+b\sqrt2$ by the
invertible transformation recorded in the source proofs. Their discriminants
are
\[
 2^{15}(b')^2(2a^2-b^2)^4(a^2-2b^2)^2.
\]
The equation $h^4=D_0$, with $D_0=(3-\sqrt2)^2$, would force a rational
fourth power to equal $49$ after taking norms, which is impossible.
The quadratic factors are anisotropic over $\Q$.
Also $b'=-4a^3+6ab^2+3b^3$ has no nonzero rational zero: for $b\ne0$
use Eisenstein at $3$ on $4X^3-6X-3$, and treat $b=0$ directly.
Thus the separability hypothesis holds at every rational retained base
point, including those with no rational elliptic witnesses.

The exact values of $(b,n,A_*,t)$ are
\[
\begin{array}{c|r|r|r|r}
 &b&n&A_*&t\\\hline
25&262144&1048576&10545336&3861794\\
27&262144&2097152&15810084&7068076\\
13\ (121/100)&64&256&523872&285114\\
13\ (16/15)&64&256&486528&273690.
\end{array}
\]
The degree sum includes respectively $254,179,3,3$ cases in which the
tail coefficients dominate the truncated determinant bound. Omitting these
cases would give an incorrect numerical bound.

The producer and a separate checker verify the exact bounds. Thirteen small
algebras are also computed using exact multiplication matrices: these
include vanishing radicands and excluded nilpotents, nonreal point and pole
roots, repeated values of $T$, a zero intermediate characteristic coefficient,
a very small radicand, an empty nonzero locus, and a base algebra with
two point generators and radicands depending on both. A deliberate omission of
$E_r$ is rejected by the empty-locus test. These are finite tests of the
formulas; the universal assertions are proved above. No large source
polynomial has been expanded, and no non-author review is represented by
these computations.

The new specifications preserve the source projections for every retained
base point. Consequently the arithmetic soundness and completeness proofs,
including their quantifier order and exceptional input treatment, apply
without change. The source construction for 13 variables is multiplied by
$X$ after conversion; this adds one to its degree bound, preserves the
coefficient multiset and support size, and introduces no variable.

\begin{thebibliography}{9}
\bibitem{Active} \emph{Separation in finite algebras on the nonzero locus},
repository companion source, \nolinkurl{papers/finite_algebra/active_finite_algebra_20261109_v15.tex}.
The present construction replaces its separating element and characteristic-polynomial bounds.
\bibitem{Source25} \emph{A defining polynomial in twenty-five variables},
\nolinkurl{papers/few_variables25/few_variables25_merged_20261110_v19.tex}.
\bibitem{Source27} \emph{An explicit defining polynomial in twenty-seven variables},
\nolinkurl{papers/few_variables27/explicit27_20261110_v19.tex}.
\bibitem{Source13} \emph{Thirteen variables from direct elliptic index comparison},
\nolinkurl{papers/conditional_few_variables/direct_input13_20261110_v22.tex}.
\bibitem{Sun} Zhi-Wei Sun, \emph{$\Z$ is Diophantine over $\Q$ with 28 unknowns},
arXiv:2611.14605v1, 10 November 2026,
\url{https://arxiv.org/html/2611.14605v1}.
Its finite-root combination is related background; no arithmetic theorem
from that paper is newly assumed in this conversion.
\end{thebibliography}
\end{document}
